\documentclass[../../main.tex]{subfiles}
\begin{document}

\begin{flushright}
\footnotesize\textit{【自動文字起こし・要確認】}
\end{flushright}

{\bf [解]}
2つの整数解を$\alpha,\beta$，2つの虚数解を$z,\bar z$とおく．明らかに$\alpha\beta z\bar z\ne0$．
\begin{align}
a&=-(\alpha+\beta+z+\bar z)\quad\cdots①\\
1&=\alpha\beta z\bar z=\alpha\beta|z|^2\quad\cdots②\\
c&=-\alpha\beta z\bar z\left(\dfrac{1}{\alpha}+\dfrac{1}{\beta}+\dfrac{1}{z}+\dfrac{1}{\bar z}\right)\\
&=-\left(\dfrac{1}{\alpha}+\dfrac{1}{\beta}+\dfrac{z+\bar z}{|z|^2}\right)\quad(\because②)\quad\cdots③\\
b&=\alpha\beta+(\alpha+\beta)(z+\bar z)+|z|^2\quad\cdots④
\end{align}

①から，$z+\bar z=2\mathrm{Re}(z)\in\mathbb Z$とかける（$\because\alpha,\beta\in\mathbb Z$）．$z+\bar z=t$（$t\in\mathbb Z$）とおく．④から同様に$|z|^2\in\mathbb Z$が従う．したがって②から，
\begin{align}
|z|^2=1,\ \alpha\beta=1\quad\therefore\ |z|=1,\ (\alpha,\beta)=(1,1),(-1,-1)\quad\cdots⑤
\end{align}
とわかる．したがって，$z$を見ると，$|z|=1$かつ$\mathrm{Re}(z)=t/2$から$\mathrm{Im}(z)^2=1-t^2/4$であり，$z$は虚数（$\mathrm{Im}(z)\ne0$）だから$|t|<2$，すなわち$t=0,\pm1$である．このとき
\begin{align}
z=\dfrac{t}{2}\pm i\sqrt{1-\dfrac{t^2}{4}}\qquad(t=0:\ z=\pm i,\quad t=1:\ z=\dfrac{1\pm\sqrt{3}i}{2},\quad t=-1:\ z=\dfrac{-1\pm\sqrt{3}i}{2})
\end{align}

\bigskip
$1^\circ$ $(\alpha,\beta)=(1,1)$の時

①④に代入して
\begin{align}
a=-(2+t),\quad b=2+2t,\quad c=-(2+t)
\end{align}
となり，
\begin{align}
f(x)=x^4-(2+t)x^3+2(1+t)x^2-(2+t)x+1=(x-1)^2(x^2-tx+1)
\end{align}

②から，$t=0,\pm1$のいずれも$x^2-tx+1=0$の判別式$t^2-4<0$を満たし十分で，
\begin{align}
(a,b,c)=(-2,2,-2),\ (-3,4,-3),\ (-1,0,-1)
\end{align}

$2^\circ$ $(\alpha,\beta)=(-1,-1)$の時
\begin{align}
a=-(-2+t)=2-t,\quad b=2-2t,\quad c=-(-2+t)=2-t
\end{align}

で，$1^\circ$と同様にしていずれも十分で
\begin{align}
(a,b,c)=(2,2,2),\ (1,0,1),\ (3,4,3)
\end{align}

以上から，複号同順として
\begin{align}
(a,b,c)=(\pm2,2,\pm2),\ (\pm1,0,\pm1),\ (\pm3,4,\pm3)．
\end{align}

\newpage

\end{document}
